\documentclass[11pt]{article}
\usepackage{amsmath,amssymb,amsthm}
\usepackage[margin=1in]{geometry}
\usepackage{hyperref}
\newtheorem{proposition}{Proposition}
\title{A distinct-level test for compact fixed-direction decoders}
\author{Research checkpoint, branch ar/027ac8030b43}
\date{November 24, 2026}
\begin{document}
\maketitle

\section{The degree is controlled by distinct input levels}

Let $X$ be a finite set of $N$ row inputs and consider the family
\[
 F(x)=\tau+\sum_{h=1}^H\frac{a_h}{b_h+t_h(x)},
 \qquad b_h+t_h(x)>0.
\]
Each map $t_h:X\to\mathbb R$ is fixed independently of the labels. Suppose
its image $T_h$ has at most $L$ elements, where $L\geq1$. The offsets,
residues, and bias may vary arbitrarily with the labels. Compactness of
the offset range is unnecessary for the following obstruction.

\begin{proposition}
Every sign pattern of this reciprocal family belongs to a family of $N$
polynomials of total degree at most $L$ in $p=2H+1$ real parameters.
If it shatters $X$ and $N\geq p$, then
\[
 2^N\leq\left(\frac{4eLN}{p}\right)^p,
 \qquad
 \frac Np\leq\log_2\left(4eL\frac Np\right).
\]
In particular, $p>N/(4\log_2(4eL))$.
\end{proposition}
\begin{proof}
For each head put
\[
 Q_h(b)=\prod_{u\in T_h}(b+u),\qquad z_h=a_h/Q_h(b_h).
\]
The denominator positivity makes $Q_h(b_h)$ nonzero. At a row $x$,
\[
 \frac{a_h}{b_h+t_h(x)}
   =z_h\prod_{u\in T_h\setminus\{t_h(x)\}}(b_h+u).
\]
The right side has total degree at most $|T_h|\leq L$ in $(z_h,b_h)$.
Treating all $z_h,b_h,\tau$ as independent parameters only enlarges the
family. Applying Warren's strict-sign bound to these $N$ polynomials
proves the displayed inequality. Write $u=N/p\geq1$ and
$A=\log_2(4eL)>3$. If $u\geq4A$, then $u>4$ and
$\log_2u\leq u/2$, so
$u\leq A+\log_2u\leq3u/4$, a contradiction.
\end{proof}

The algebraic cancellation, polynomial evaluation identity, and degree
bound are kernel checked in
\texttt{formalization/HeadComplexity/Results/LargeDomainFiniteLevels.lean}:
\begin{itemize}
\item \texttt{HeadComplexity.LargeDomain.reciprocal\_eq\_level\_product};
\item \texttt{HeadComplexity.LargeDomain.finiteLevelScorePolynomial\_eval};
\item \texttt{HeadComplexity.LargeDomain.finiteLevelScorePolynomial\_totalDegree\_le}.
\end{itemize}
The Lean formulation permits a list of $L$ levels with repetitions,
which also permits padding smaller level sets to a common bound.
Warren's theorem and the analytic counting consequences in this note
are written proofs using that external theorem; they are not represented
as unconditional Lean declarations here. No new bound on the original
$W(n)$ is claimed. The strict-sign form used above is Theorem 2.1 of
Goldberg and Jerrum, \emph{Bounding the Vapnik--Chervonenkis dimension
of concept classes parameterized by real numbers}, available at
\url{https://www.cs.ox.ac.uk/people/paul.goldberg/papers/goldbergjerrum.pdf};
it attributes the theorem to Warren's 1968 paper,
\href{https://doi.org/10.1090/S0002-9947-1968-0226281-1}{\emph{Lower bounds
for approximation by nonlinear manifolds}}.

\section{Fixed affine row numerators}

The conclusion survives the broader family
\[
 F(x)=\tau+\sum_{h=1}^H\frac{A_h(x)+a_h}{b_h+t_h(x)},
\]
where each $A_h$ is fixed independently of the labels. It is irrelevant
for this count whether $A_h$ is affine: its value at each row is a fixed
coefficient. Put $u_h=1/Q_h(b_h)$ and $v_h=a_h/Q_h(b_h)$. Then
\[
 \frac{A_h(x)+a_h}{b_h+t_h(x)}
  =(A_h(x)u_h+v_h)
       \prod_{t\in T_h\setminus\{t_h(x)\}}(b_h+t).
\]
Allowing $u_h,v_h,b_h$ to be independent gives degree at most $L$ in
$p=3H+1$ parameters. The same Warren inequality and lower bound on $p$
follow. The restrictions $u_hQ_h(b_h)=1$ and
$v_h=a_hu_h$ need not be counted because dropping them enlarges the
polynomial family.

This argument does not apply with $O(H)$ parameters when the $k$
affine row coefficients of each $A_h$ may themselves vary independently
with the labels. That case has $\Theta(Hk)$ numerator parameters. Nor
may the fixed directions $t_h$ be silently allowed to depend on each
independent column's labels.

\section{The Johnson and Hamming-layer candidate fails the test}

Take $X=\{0,1\}^k$ and one head for each central subset
$S\subseteq[k]$, $|S|=\lfloor k/2\rfloor$, with
\[
 t_S(x)=T|x|-\epsilon|x\cap S|,\qquad T>\epsilon>0.
\]
All coordinate slopes are strictly positive. There are at most
$(k+1)^2$ possible pairs $(|x|,|x\cap S|)$, hence at most that many
levels for each direction. The proposition gives
\[
 H=\Omega\left(\frac{2^k}{\log k}\right)
\]
for universal sign decoding, even with fixed affine row numerators.
Thus the apparent width-scale budget
$H=\binom{k}{\lfloor k/2\rfloor}=\Theta(2^k/\sqrt{k})$
cannot shatter the entire cube in this family for all large $k$.
Exact interpolation on a single central layer does not overcome the
independently chosen labels on all other layers.

A finite numerical experiment at $k=4$, $H=6$, $T=1$,
$\epsilon=0.01$, and offsets in $(1,9/8)$ observed $30437$ of the $32768$
sign complement classes after $400$ offset samples. The arrangement
calculation was floating point, not an exact certificate. Its missing
patterns are not a finite impossibility proof. The asymptotic obstruction
above is independent of this computation.

More generally, if every head's coordinate weights take values in a
fixed set of $r$ numbers, then $L\leq(k+1)^r$. A bounded number of
weight classes therefore has the same obstruction. Integer weights in
$\{1,\ldots,B\}$ give $L\leq kB+1$.

\section{Required level growth and the next route}

Suppose $H\leq C2^k/G(k)$ with $G(k)\to\infty$ and $G(k)=o(2^k)$.
The exact inequality above implies
\[
 \log_2L\geq \frac{2^k}{p}
             -\log_2\left(\frac{2^k}{p}\right)-\log_2(4e).
\]
Consequently
\[
 \liminf_{k\to\infty}\frac{\log_2L}{G(k)}\geq\frac1{2C}
\]
for constant row numerators, and the conservative constant is $1/(3C)$
for fixed affine row numerators. In particular, a saving $G(k)=k^\delta$
requires $L\geq\exp(\Omega(k^\delta))$. Improving the current
$H=O(2^k/\log k)$ scale by an unbounded factor $g(k)$ requires
$L\geq k^{\Omega(g(k))}$.

The next construction must therefore use exponentially many distinct
subset-sum levels for a fixed-power gain. Generic unequal positive
weights meet this necessary condition. Arbitrarily small perturbations
can already make all $2^k$ levels distinct, so the present theorem does
not rule out perturbed Johnson or uniform-layer families. Replacing
nearby levels by equal ones is unjustified when residues and precision
are unrestricted. A successful perturbation argument must control those
large coefficients and the strict margins explicitly.

Merely making all levels distinct by choosing a common binary-weight
linear form is also insufficient in the constant-numerator subfamily:
if every $t_h$ is a scalar multiple of one common form, the score is a
univariate rational function with numerator degree at most $H$. Alternating
labels on $N$ distinct ordered levels require at least $N-1$ real zeros,
so $H\geq N-1$. The remaining opportunity needs many distinct levels
and enough nonparallel directions to avoid this one-dimensional bound.
\end{document}
